\section{Security}

\autoref{tab:security} lists the threats I considered and how the platform answers each.

\begin{table}[H]
	\centering
	\begin{tabular}{|p{0.28\textwidth}|p{0.62\textwidth}|}
		\hline
		\textbf{Threat} & \textbf{Measure} \\
		\hline
		Stolen password database & Passwords are stored only as bcrypt hashes with a salt per password~\cite{owasp}. \\
		\hline
		Session theft by injected scripts & The session token lives in an HTTP-only cookie that page scripts cannot read; React escapes all rendered text. \\
		\hline
		Forged requests from other sites & The cookie uses \texttt{SameSite=Lax}, and only the dashboard's own origin is allowed. \\
		\hline
		Access to other users' data & Every match query filters by the owner; other users' matches answer ``not found'', so their IDs cannot be probed. \\
		\hline
		Malicious uploads & Extension, type and the file's first bytes are checked; files get random names and a size limit. \\
		\hline
		Paths outside the storage & Every stored path is resolved and must stay inside the storage folder. \\
		\hline
		Invalid input & All request bodies and settings are validated with Zod. \\
		\hline
		Leaked secrets & Secrets live only in \texttt{.env}, which Git ignores; the language model receives statistics only. \\
		\hline
	\end{tabular}
	\caption{Security measures}
	\label{tab:security}
\end{table}

\section{Implementation decisions}

The code was typed throughout: TypeScript for the API and dashboard, type hints and dataclasses for the worker, and a database client which was automatically generated by the schema. When the output format of the results changed, the compiler would flag all places in the dashboard that were reading them.

Because I stream the data (i.e., video frames are pulled into memory in batch sizes, uploads directly to disk, the GPU upload is done piece-by-piece and the video player loads with range requests), memory remained relatively stable. A full half peaked at 3.6--4.4\,GB of RAM. I use Pandas and NumPy operations to process tracking data, as opposed to using Python loops, and I operate on whole columns. A half contains approximately 840,000 detections. All of the steps after pitch mapping combined took less than 3 minutes.

Each threshold in the pipeline has a name constant defined in \texttt{config.py} and a comment describing its value and when it was measured and on what data, therefore, if I want to run another experiment, I simply need to change one line. Comments describe why something is happening. All generic code is shared: there is one drawing function that can be used to draw both the board and player tab, there is one statistics function that can work either per track or per position, and the same \texttt{run\_\allowbreak{}analysis()} function will work whether a local or a GPU analysis is being performed.

There is an expectation of errors occurring (i.e., jobs are retried, results are written atomically, values that cannot be stored in JSON are replaced before writing, the GPU will fall-back to local analysis, and if a job is cancelled then no partially-written files will remain).
